\documentclass[12pt]{article}

\usepackage{comment}
\usepackage{amsmath}
\usepackage{amssymb}  % for \nmid

\newcommand{\D}{{\mathbb D}}
\newcommand{\G}{{\mathbb G}}
\newcommand{\Z}{{\mathbb Z}}
\newcommand{\N}{{\mathbb N}}
\newcommand{\Q}{{\mathbb Q}}
\newcommand{\R}{{\mathbb R}}
\newcommand{\succc}{{\rm succ}}
\newcommand{\pred}{{\rm pred}}

\newcommand{\lc}{\left\lceil}
\newcommand{\rc}{\right\rceil}
\newcommand{\Ceil}[1]{\left\lceil {#1}\right\rceil}
\newcommand{\ceil}[1]{\left\lceil {#1}\right\rceil}
\newcommand{\floor}[1]{\left\lfloor{#1}\right\rfloor}

\begin{document}


\centerline{Homework 07, MORALLY Due April 7} 

\newif{\ifshowsoln}
\showsolntrue % comment out to NOT show solution inside \ifshowsoln and \fi blocks.

\begin{enumerate}
\item
(0 points)
What is your name.
\textbf{Gurkeerat Singh}

\vfill 

\centerline{\bf GO TO THE NEXT PAGE}

\newpage


\item
(50 points)
\begin{enumerate}
\item
(0 points) 
(You may want to write a program to get data.)

$(a,b,c)$ means the NIM game where you can remove $a$ or $b$ or $c$ stones.

Below the condition XXX will be of the form

\bigskip

$n\equiv a_1,a_2,\ldots,a_k \pmod m.$

\bigskip

\item 
(7 points) 
Fill in the XXX:

{\it In the game $(1,3,5)$, Player II wins when the game begins with $n$ stones iff  $XXX(n)$.} 
\begin{itemize}
    \item \textbf{Solution: Player II wins when the game begins with $n$ stones iff $n\equiv 0,2 \pmod 8.$}
\end{itemize}

\item
(7 points) 
Fill in the XXX:

{\it In the game $(1,4,6)$, Player II wins when the game begins with $n$ stones iff  $XXX(n)$.} 
\begin{itemize}
    \item \textbf{Solution: Player II wins when the game begins with $n$ stones iff $n\equiv 0,2 \pmod 5.$}
\end{itemize}

\item
(7 points) 
Fill in the XXX:

{\it In the game $(1,5,7)$, Player II wins when the game begins with $n$ stones iff  $XXX(n)$.} 
\begin{itemize}
    \item \textbf{Solution: Player II wins when the game begins with $n$ stones iff $n\equiv 0 \pmod 2.$}
\end{itemize}

\item
(7 points) 
Fill in the XXX:

{\it In the game $(1,6,8)$, Player II wins when the game begins with $n$ stones iff  $XXX(n)$.}
\begin{itemize}
    \item \textbf{Solution: Player II wins when the game begins with $n$ stones iff $n\equiv 0,2,4 \pmod 7.$}
\end{itemize} 

\item
(22 points) 
Fill in the XXX

{\it In the game $(1,m,m+2)$, Player II wins when the game begins with $n$ stones iff  $XXX(n,m)$.} 
\begin{itemize}
    \item \textbf{Solution: Player II wins when the game begins with $n$ stones iff $n\equiv 0,2,4,\ldots,2(m-1) \pmod{2m+1}.$}
\end{itemize}

\end{enumerate} 

\vfill

\centerline{\bf GO TO THE NEXT PAGE}

\newpage

\item 
(50 points) 
Let $a\in\N$.  Consider the recurrence

$T_a(0)=10$

$T_a(n) = 2T(\floor{\frac{n}{2}}) + an$.

\begin{enumerate}
\item
(0 points) 
Write a program that will, given $a$, output

$T_a(0)$, $T_a(1)$, $\ldots$, $T_a(1000)$. 

\item 
(0 points) 
Run the program for $1\le a\le 20$. 

\item
(50 points) Using your data make a conjecture along the lines of:

$T_a(n)$ is roughly $XXX(a,n)$

(For example, $T_a(n) = a^2n^3$ which is NOT the answer.) 
\begin{itemize}
    \item \textbf{Solution: $T_a(n)$ is roughly $an\log n$.}
\end{itemize}
\end{enumerate}


\end{enumerate} 

\end{document}
